%%%%%%%%%%%%%%%%%%%%%%%%%%%%%%%%%%%%%%%%%%%%%%%%%%%%%%%%%%%%%%%%%%%%%%
% Lớp 4 | Chủ đề 1. Ôn tập và bổ sung | Giải bài toán có ba bước tính
% Dạng 4D1?5-1: Bài toán ba bước tính (mua bán, chia đều)
%%%%%%%%%%%%%%%%%%%%%%%%%%%%%%%%%%%%%%%%%%%%%%%%%%%%%%%%%%%%%%%%%%%%%%


% nguồn: 2026-11-18_lop4_chu-de-1-tu-luan (nhập 18/11/2026)
\begin{bt}%[4D1N5-1]%[ma:564aaa62]
% nguon: Toán tư duy lớp 4 – Bài 5 SGK Toán 4 KNTT
	Một cửa hàng nhận về $4$ thùng nước mắm, mỗi thùng $24$ chai. Buổi sáng cửa hàng bán $30$ chai, buổi chiều bán $28$ chai. Hỏi cửa hàng còn lại bao nhiêu chai nước mắm?
	\loigiai{
		Bước 1. Số chai nhận về: $24 \times 4 = 96$ (chai).\\
		Bước 2. Số chai đã bán: $30 + 28 = 58$ (chai).\\
		Bước 3. Số chai còn lại: $96 - 58 = 38$ (chai).\\
		Đáp số: $38$ chai.
	}
\end{bt}

% nguồn: 2026-11-18_lop4_chu-de-1-tu-luan (nhập 18/11/2026)
\begin{bt}%[4D1H5-1]%[ma:85bd694a]
% nguon: Toán tư duy lớp 4 – Bài 5 SGK Toán 4 KNTT
	Hà mua $3$ quyển vở, mỗi quyển giá $8\,000$ đồng và $2$ cái bút, mỗi cái giá $6\,500$ đồng. Hà đưa người bán tờ $50\,000$ đồng. Hỏi người bán phải trả lại Hà bao nhiêu tiền?
	\loigiai{
		Bước 1. Tiền mua vở: $8\,000 \times 3 = 24\,000$ (đồng).\\
		Bước 2. Tiền mua bút: $6\,500 \times 2 = 13\,000$ (đồng).\\
		Bước 3. Tiền trả lại: $50\,000 - (24\,000 + 13\,000) = 13\,000$ (đồng).\\
		Đáp số: $13\,000$ đồng.
	}
\end{bt}

% nguồn: 2026-11-18_lop4_chu-de-1-tu-luan (nhập 18/11/2026)
\begin{bt}%[4D1V5-1]%[ma:48f62644]
% nguon: Toán tư duy lớp 4 – Bài 5 SGK Toán 4 KNTT
	Một trường có khối Ba, khối Bốn và khối Năm. Khối Ba có $156$ học sinh, khối Bốn nhiều hơn khối Ba $24$ học sinh, khối Năm ít hơn khối Bốn $18$ học sinh. Hỏi ba khối có tất cả bao nhiêu học sinh?
	\loigiai{
		Khối Bốn: $156 + 24 = 180$ (học sinh).\\
		Khối Năm: $180 - 18 = 162$ (học sinh).\\
		Cả ba khối: $156 + 180 + 162 = 498$ (học sinh).\\
		Đáp số: $498$ học sinh.
	}
\end{bt}

% nguồn: 2026-11-18_lop4_chu-de-1-tu-luan-bo-sung (nhập 18/11/2026)
\begin{bt}%[4D1H5-1]%[ma:f992f3e5]
% nguon: Toán tư duy lớp 4 – Bài 5 SGK Toán 4 KNTT
	Một xe tải chở gạo $3$ chuyến, mỗi chuyến $1\,250$ kg. Số gạo đó được chia đều cho $25$ cửa hàng. Hỏi mỗi cửa hàng nhận được bao nhiêu ki-lô-gam gạo?
	\loigiai{
		Bước 1. Số gạo xe chở được: $1\,250 \times 3 = 3\,750$ (kg).\\
		Bước 2. Mỗi cửa hàng nhận: $3\,750 : 25 = 150$ (kg).\\
		Đáp số: $150$ kg gạo.
	}
\end{bt}
